Dans l'infrarouge, la peau humaine se comporte presque comme un corps noir. On considère une personne dont la peau est à \qty{\v{Tc}}{\degreeCelsius} et dont la surface corporelle est de \qty{\v{S}}{m^2}.

\begin{enumerate}[label=\alph*)]
    \item À quelle longueur d'onde le rayonnement émis par la peau est-il maximal ? Pourquoi les caméras thermiques détectent-elles cette région du spectre ?
    \item Déterminez la puissance émise par mètre carré de peau.
    \item Déterminez la puissance totale émise par la personne. Pourquoi ne perd-elle pas réellement autant d'énergie ?
\end{enumerate}
